\section{手拉脚与中点}
\begin{problem}{题 1 共用一个底角顶点与一个直角顶点}
如图，在 \(\triangle ADE\) 和 \(\triangle ABC\) 中，\(AD=AE\)，\(AB=BC\)，\(\angle DAE=\angle ABC=90^\circ\)。延长 \(DB\) 到 \(F\)，使 \(BF=DB\)，连接 \(CF\)、\(CE\)。求证：\(CF=CE\)，\(CF\perp CE\)。

各点方向按图取定。尝试两种构造：延长 \(AB\)；或者把 \(AC\) 绕 \(A\) 旋转 \(90^\circ\)。

\diagram{\hhOne}

\hint{\(B\) 已经是 \(DF\) 的中点，能否让它也成为一条新线段的中点？}

\end{problem}
\clearpage
\section{脚拉脚与两次倍长}
\begin{problem}{题 2 两个底角顶点重合}
如图，\(\triangle ABC\) 和 \(\triangle DCE\) 都是等腰直角三角形，\(\angle BAC=\angle CDE=90^\circ\)，即 \(AB=AC\)、\(DC=DE\)。\(D\) 在 \(BA\) 向 \(A\) 外的延长线上，\(E\) 与 \(C\) 位于直线 \(AB\) 的两侧。连接 \(BE\)，求证：\(BE\perp BC\)。

\diagram{\hhTwo}

\hint{分别倍长 \(ED\)、\(BA\)，使公共点 \(C\) 成为两个新等腰直角三角形的直角顶点。}

\end{problem}
\clearpage
